\documentclass[10pt,a4paper]{amsart}
\usepackage[margin=19mm]{geometry}
\usepackage{amsmath,amssymb,mathtools}
\usepackage[scr=rsfs,cal=euler]{mathalfa}
\usepackage{xfrac}
\usepackage[unicode,hidelinks,bookmarks=false]{hyperref}
\newcommand{\ie}{i.e.\ }
\newcommand{\cf}{cf.\ }
\newcommand{\resp}{respectively\ }
\newcommand{\Subsection}{\subsection}
\providecommand{\nobreakdash}{\nobreak\hbox{-}}
\setlength{\emergencystretch}{2em}
\multlinegap0pt
\usepackage{amssymb}
\usepackage{enumerate}
\newtheorem{lemma}{Lemma}
\newcommand{\cC}{\mathcal C}
\newcommand{\cX}{\mathcal X}
\title{Upper semicontinuity of algebraic dimension}
\author{Sheng Rao}
\date{Source: arXiv:2407.02022v3 (CC BY 4.0)}
\begin{document}
\maketitle

\noindent\emph{Source and licence.} Upper semicontinuity of algebraic dimension. Version arXiv:2407.02022v3 (CC BY 4.0), distributed by arXiv under the Creative Commons Attribution 4.0 International licence, http://creativecommons.org/licenses/by/4.0/, as recorded in the arXiv licence metadata for this exact version and shown on the abstract page https://arxiv.org/abs/2407.02022v3. Full licence notice: \texttt{licenses/arxiv-2407.02022-CC-BY-4.0.txt}.\par\smallskip

\noindent\textit{Editorial scope.} Counted unit: the article's lemma lem:revised-global-component (source0.tex lines 1105--1123). The article's complete original proof of it is delivered with it (source0.tex lines 1125--1156, one \texttt{proof} environment of 32 lines). No mathematics is written, repaired or re-derived: the statement and the proof are the article's own lines, in the article's own notation, and no retained line is substituted or re-flowed.  No further source-local context is needed: every cross-reference of every retained line resolves inside the two delivered spans.  Nothing was omitted from any retained span, so the package carries no omission record.  No retained line contains a citation, so the wrapper carries no bibliography and no bibliography entry.  No retained line contains a figure environment, an image-inclusion command or drawing code, so no image is delivered and the package declares no asset dependency.  Editorial formatting only, in addition: an AMS \texttt{amsart} preamble that restates the article's own declarations and macros used by the retained lines, namely the environments \texttt{lemma} on separate counters, together with the standard packages those lines need.  The retained environments print in this document as Lemma 1 in order of appearance, whereas the article prints the same items with the numbering of its own full document.  The complete native source of the article as distributed in the arXiv e-print for this exact version is bundled unchanged as \texttt{sources/161679/source0.tex}.
\begin{lemma}\label{lem:revised-global-component}
Let $Y$ be a reduced complex space, let $Y^\circ\subset Y$ be open, and let
$A^\circ$ be an irreducible component of $Y^\circ$.  Then there is an irreducible component $A$ of $Y$ such that
\begin{equation}\label{eq:revised-component-restriction}
 A^\circ\subset A\cap Y^\circ.
\end{equation}
Moreover, $A^\circ$ is an irreducible component of $A\cap Y^\circ$,
$\dim A=\dim A^\circ$, and $A^\circ$ contains a nonempty subset which is
open in $A$.  In particular, with
\[
 \mathscr C:=\cC_{n-1}(\cX/\Delta),\qquad
 \mu:\mathscr C\to\Delta,
\]
there is an irreducible component $\Gamma\subset\mathscr C$ satisfying
\begin{equation*}
 \Gamma_U\subset\Gamma\cap\mu^{-1}(U),
\end{equation*}
and $\Gamma_U$ contains a nonempty open subset of $\Gamma$.
\end{lemma}

\begin{proof}
Complex spaces are locally Noetherian, and their local irreducible
components are locally finite.  In a neighbourhood meeting only finitely
many local components, the intersections of the local branches of
$A^\circ$ with the other components are proper analytic subsets.  We may
therefore choose a point $a\in A^\circ$ which lies on no other local
component of $Y^\circ$ and at which the germ of $A^\circ$ is irreducible.

A local irreducible component through $a$ is
contained in an irreducible component $A$ of $Y$.  Hence
$A\cap A^\circ$ is a closed analytic subset of the irreducible space
$A^\circ$ which contains a nonempty open subset.  The identity principle
gives $A^\circ\subset A$.

If an irreducible closed analytic subset of
$A\cap Y^\circ$ contains $A^\circ$, then it is also an irreducible closed
analytic subset of $Y^\circ$; maximality of $A^\circ$ forces equality.
Thus $A^\circ$ is an irreducible component of $A\cap Y^\circ$.

The other irreducible components of $A\cap Y^\circ$ form a locally finite
family, and their intersections with $A^\circ$ are proper analytic
subsets.  Choose a point of $A^\circ$ outside their union and shrink a
neighbourhood in $A\cap Y^\circ$.  The resulting nonempty subset is open
in $A$ and is contained in $A^\circ$.  In particular,
$\dim A=\dim A^\circ$.

One cannot in general strengthen \eqref{eq:revised-component-restriction} to
$A\cap Y^\circ=A^\circ$: for the irreducible curve $t=y^2$, the inverse
image of a small simply connected disc away from $t=0$ has two components.
Apply the valid containment statement to $Y=\mathscr C$ and
$Y^\circ=\mu^{-1}(U)$.
\end{proof}

\end{document}
